%قبل از شروع مسئله ترجیحاً چیزی ننویسید. اما اگر خواستید بنویسید، در قالب لیستی (فقرات و شمارش) بنویسید.
%برای اضافه کردن نمره، بعدی از شروع‌مسئله، [x] را اضافه کنید و به جای x نمره مورد نظر را قرار دهید. برای زیرمسئله‌ها هم به همین صورت است.
%\شروع‌مسئله[10]
\شروع‌مسئله [15]
%\متن‌سیاه{سوال عملی اول:}\\
%\متن‌سیاه{دراپ‌تست}\\

%\متن‌سیاه{ورودی:}\\
%در خط اول ورودی، ابتدا k و سپس m داده می‌شود.\\
%\متن‌سیاه{خروجی:}\\
%در تنها سطر خروجی، تعداد آیفون‌های مورد نیاز می‌آید. اگر در هیچ صورتی چنین کاری ممکن نبود ۱- چاپ می‌شود.\\
%\متن‌سیاه{محدودیت‌ها:}\\
%$1\leq k\leq 10^4$

%$m\leq10^9$
%\\
%\textbf{\متن‌سیاه{ورودی و خروجی نمونه}\\
%% تست اول
%\begin{latin}
% \begin{tabularx}{\textwidth}{|X | X|}
% \hline
% \textbf{stdin} &
% \textbf{stdout}\\
% \hline
% 100 100&
% 1\\
% \hline
% 
% \end{tabularx}
%\end{latin}
%
%% تست دوم
%
%\begin{latin}
% \begin{tabularx}{\textwidth}{|X | X|}
% \hline
% \textbf{stdin} &
% \textbf{stdout}\\
% \hline
% 340 20&
% 3\\
% \hline
% 
% \end{tabularx}
%\end{latin}}
% صورت اصلی مسئله را در این‌جا بنویسید.
گراف جهتدار و وزن‌دار
$G(V,E)$
داده شده است. وزن‌ها ممکن است منفی هم باشند. $s$ و $t$ دو راس گراف هستند و $k$ عددی طبیعی است.

\شروع{زیرمسائل}

\زیرمسئله[5] % مسئلۀ ۱آ
الگوریتمی از مرتبۀ زمانی 
$O(|V|\log |V| + |E|)$
ارائه دهید که مسیری از $s$ به $t$ که مینیمم یال‌هایش بیشینه است، را پیدا کند.

\زیرمسئله[5] % مسئلۀ ۱آ
الگوریتمی از  مرتبۀ 
$O(k(|V| + |E|))$
ارائه دهید که میان تمام مسیرهای $s$ به $t$ که دقیقا $k$ یال دارند، طول کوتاه‌ترین را پیدا کند.

\متن‌سیاه{راهنمایی:}
یک گراف جدید تشکیل دهید و به ازای هر راس $v$ در گراف اصلی، 
$k + 1$
راس $(v,i)$ در نظر بگیرید
$(0 \le i \le k)$
،
و یال‌ها را طوری قرار دهید که فاصلۀ راس
$(v,i)$
از $(s,0)$ برابر با $i$ باشد. سپس ثابت کنید این گراف یک \متن‌لاتین{DAG} است و کوتاه‌ترین مسیر از $(s,0)$ به $(t,k)$ را بیابید.
\زیرمسئله[5] % مسئلۀ ۱آ
فرض کنید وزن‌ها مثبت هستند.
بخش قبلی را با این تفاوت که تعداد یال‌های مسیرها \underline{حداقل} $k$ باشد، حل کنید. الگوریتم‌تان باید از مرتبۀ 
$O(|V|\log |V| + k|V| + k|E|)$
باشد.

\پایان{زیرمسائل}


\solution{
	در قسمت آ، یک گراف جهتدار داریم که در آن باید مسیری از $s$ به $t$ را 
	پیدا کنیم که کمترین یال آن مسیر، بیشینه است.
	الگوریتمی که ارائه می‌دهیم همان الگوریتم دایسترا با کمی تغییر است. 
	در الکوریتم دایسترا مسیر بهینه مسیری است که مجموع وزن یال‌هایش کمینه شود. اینجا ما نیاز داریم مسیری را پیدا کنیم که کوچک‌ترین یالش بیشینه است.
	در هر مرحله از الگوریتم دایسترا راس‌ها دو دسته‌اند. دستۀ اول شامل راس‌هایی است که کوتاهترین مسیر برای‌شان به‌دست آمده است. دسته دوم راس‌هایی هستند که کوتاهترین مسیر برای آن‌ها هنوز به‌دست نیامده و در یک هرم کمینه مرتب شده‌اند. در هر مرحله کلید هر راس در این درخت کمینه برابر با کوتاه‌ترین مسیر از $s$ به آن راس است که راس‌های میانی آن همه در دستۀ اول قرار دارند.
	دایسترا با $d[s] = 0$ و $d[v] = \infty$ برای راس‌های دیگر آغاز می‌شود.
	حالا الگوریتم دایسترا را بدین صورت تغییر می‌دهیم. 
	در ابتدا قرار دهید
	$d[s] = \infty$
	و
	$d[v] = 0$
	برای راس‌های دیگر.
	به جای هرم کمینه یک هرم بیشینه درنظر بگیرید که راس $i$ با کلید $d[i]$ در آن قرار دارد. بنابراین در ابتدا خود $s$ در راس هرم قرار دارد.
	مشابه الگوریتم دایسترا در هر مرحله راس بالای هرم ($u$) را خارج می‌کنیم و در دستۀ اول همراه با $s$ قرار می‌دهیم.
	هر یال
	$(u,v)$ 
	که از $u$ خارج می‌شود و به $v$ (یک راس در هرم) وصل است، را ریلکس می‌کنیم. یعنی قرار می‌دهیم:
	$$d[v] = \max (d[v], \min (d[u], w(u, v)))$$
	به راحتی می‌توان ثابت کرد که در هر مرحله بهترین مسیر برای راس بالای هرم به دست آمده است.
	اگر از هرم فیبوناچی استفاده کنیم، زمان اجرا مشابه دایسترای عادی برابر با
	$O(|V| \log |V| + |E|)$
	خواهد شد.
	
	در قسمت ب، یک گراف جهتدار وزن‌دار داریم و می‌خواهیم کوتاهترین مسیر از راس $s$ به راس $t$ را بیابیم که طولش دقیقا برابر با $k$ است.
	
	یک گراف جدید تشکیل دهید و به ازای هر راس $u$ در گراف اصلی، و هر
	$0\le i \le k$
	یک راس
	$(u, i)$
	بگذارید.
	به ازای هر یال از
	$u$
	به
	$v$
	در گراف اصلی و برای هر $i$، از 
	$(u,i)$
	به 
	$(v,i+1)$
	در گراف جدید یالی با همان وزن بگذارید.
	در این صورت، باید کوتاه‌ترین مسیر از 
	$(s,0)$
	به
	$(t, k)$
	را پیدا کنیم.
	اما توجه کنید که گراف جدید یک 
	DAG
	با
	$O(k|V|)$
	راس و 
	$O(k|E|)$
	یال است و پیدا کردن کوتاه‌ترین مسیر در این گراف 
	$O(k|V| + k|E|)$ 
	طول می‌کشد.
	
	در قسمت پ، باید کوتاه‌ترین مسیر از $s$ به $t$ که طولش حداقل 
	$k$
	است را پیدا کنیم.
	این کوتاه‌ترین مسیر را در نظر بگیرید. 
	$k$
	امین راس میانی این مسیر را در نظر بگیرید. آن را $v$ بنامید. مسیر از $s$ به $v$ باید کوتاه‌ترین مسیر از $s$ به $v$ به طول دقیقا $k$ باشد. ادامۀ مسیر از $v$ به $t$ نیز باید کوتاه‌ترین مسیر (بدون محدودیتی روی طول) از $v$ به $t$ باشد.
	پس می‌توانیم این‌گونه عمل کنیم. برای هر راس کوتاه‌ترین مسیر به طول $k$ از $s$ به آن را مشابه قسمت قبل به دست بیاوریم. سپس با استفاده از یک دایسترای برعکس (جهت یال ها وارونه) کوتاه‌ترین مسیر از هر راس به $t$ را به‌دست بیاوریم. در آخر تمام راس‌ها را بررسی کرده و طول دو مسیری که برای‌شان بدست آوردیم را جمع کنیم. کوچک‌ترین این مقادیر جواب مسئله است.

}

%%%%%%%%%%%%%%%%%%%%%%%%%%%%%%%%%%%%%%%%%%%%%%%%%%%%%
% برخی دستورات مفید برای تمرین‌ها. %
%%%%%%%%%%%%%%%%%%%%%%%%%%%%%%%%%%%%%%%%%%%%%%%%%%%%%

% فرمول‌های نمونه:
%$f(x) = \Theta(x)$
%$T(x, y) \leq \log(x) + 2^y + \binom{2n}{n}$
% {\tt code\_function}


% لیست موارد بدون شماره‌گذاری:
%\شروع{فقرات}
% \فقره مورد اول 
% \شروع{فقرات}
% \فقره مورد اول داخلی 
% \شروع{فقرات}
% \فقره اولین مورد خیلی داخلی 
% \فقره دومین مورد خیلی داخلی 
% \پایان{فقرات}
% \فقره مورد دوم داخلی 
% \پایان{فقرات}
% \فقره مورد دوم 
%\پایان{فقرات}

% لیست موارد شماره‌گذاری شده:
%\شروع{شمارش}
% \فقره یک مورد 
% \فقره یک مورد دیگر 
% \فقره باز هم یک مورد دیگر
%\پایان{شمارش}

% برای نوشتن یک فرمول چند خطی که یک شماره بخوره. برای حالت‌های دیگر به سایت زیر مراجعه کنید:
% https://www.overleaf.com/learn/latex/Aligning_equations_with_amsmath
%\begin{align} 
% \begin{split}
% (x+y)^3 &= (x+y)^2(x+y) \\
% &= (x^2+2xy+y^2)(x+y) \\
% &= (x^3+2x^2y+xy^2) + (x^2y+2xy^2+y^3) \\
% &= x^3+3x^2y+3xy^2+y^3
% \end{split} 
%\end{align}

% نوشتن ماتریس و آرایه دوبعدی
%\begin{align}
% A = 
% \begin{bmatrix}
% A_{11} & A_{21} \\
% A_{21} & A_{22}
% \end{bmatrix}
%\end{align}

% اضافه کردن تصویر با شرح. برای اضافه کردن تصویر بدون شرح، فقط دستور درج‌تصویر را استفاده کنید.
%\شروع{شکل}[h!]
%\تنظیم‌ازوسط
% \درج‌تصویر[پهنا=0.5\پهنای‌متن]{img.jpg}
%\شرح{شرح تصویر}
%\برچسب{تصویر:۱}
%\پایان{شکل}
%\پایان{زیرمسائل}

%کد زیر را ببینید.
%%\lstinputlisting{source_filename.py}
%\begin{latin}
%\begin{lstlisting}
%# Import the modules
%import sys
%import random
%
%ans = True
%
%while ans:
%question = raw_input("Ask the magic 8 ball a question: (press enter to quit) ")
%
%answers = random.randint(1,8)
%
%if question == "":
%sys.exit()
%
%elif answers == 1:
%print "It is certain"
%
%elif answers == 2:
%print "Outlook good"
%
%elif answers == 3:
%print "You may rely on it"
%
%elif answers == 4:
%print "Ask again later"
%
%elif answers == 5:
%print "Concentrate and ask again"
%
%elif answers == 6:
%print "Reply hazy, try again"
%
%elif answers == 7:
%print "My reply is no"
%
%elif answers == 8:
%print "My sources say no"
%\end{lstlisting}
%\end{latin}